\documentclass[11pt]{article}

% ---------- packages ----------
\usepackage[margin=1in]{geometry}
\usepackage{amsmath, amssymb, amsthm}
\usepackage{mathtools}
\usepackage{empheq}
\usepackage{bm}
\usepackage{xcolor}
\usepackage{hyperref}
\hypersetup{colorlinks=true, linkcolor=blue, urlcolor=blue, citecolor=blue}

% ---------- notation shortcuts ----------
\newcommand{\R}{\mathbb{R}}
\newcommand{\E}{\mathbb{E}}
\newcommand{\KL}{\mathrm{KL}}
\newcommand{\D}{\mathcal{D}}
\newcommand{\given}{\,\vert\,}
\newcommand{\pol}{\pi_{\theta}}          % policy being trained
\newcommand{\piref}{\pi_{\mathrm{ref}}}  % frozen reference policy
\newcommand{\polstar}{\pi^{\star}}       % optimal policy
\newcommand{\yw}{y_w}                     % chosen / winning response
\newcommand{\yl}{y_l}                     % rejected / losing response
\DeclareMathOperator{\sigmoid}{\sigma}
\newtheorem{lemma}{Lemma}
\newtheorem{proposition}{Proposition}

\title{Direct Preference Optimization: Derivation Notes}
\author{dpo-qwen}
\date{\today}

\begin{document}
\maketitle

\begin{abstract}
Working notes deriving the DPO objective from the RLHF objective, so the
implementation that follows is grounded in the math rather than copied from a
library. The throughline: \emph{the RLHF reward-maximization problem has a
closed-form optimal policy, which lets us rewrite the reward in terms of the
policy itself, which collapses the whole RL pipeline into a single
classification-style loss on preference pairs.}
\end{abstract}

\tableofcontents

% ============================================================
\section{Setup and notation}
% ============================================================

We have a prompt $x$ and a response $y$. A policy $\pi(y \given x)$ is an
autoregressive language model: a distribution over response sequences given a
prompt.

\begin{itemize}
  \item $\piref(y \given x)$ — the \textbf{reference policy}: a frozen copy of the
        model after supervised fine-tuning (SFT). It is our anchor.
  \item $\pol(y \given x)$ — the \textbf{policy} we are optimizing, initialized
        from $\piref$.
  \item $\D = \{(x, \yw, \yl)\}$ — a preference dataset. For each prompt $x$, a
        human (or model) judged response $\yw$ as preferred over $\yl$.
  \item $r(x, y)$ — a (latent) scalar reward measuring how good $y$ is for $x$.
  \item $\beta > 0$ — a temperature controlling how far $\pol$ may drift from
        $\piref$.
\end{itemize}

% ============================================================
\section{The RLHF objective}
% ============================================================

Classic RLHF maximizes expected reward while staying close to the reference
policy via a KL penalty:
\begin{equation}
  \max_{\pol} \;
  \E_{x \sim \D} \,
  \E_{y \sim \pol(\cdot \given x)}
  \Big[\, r(x, y) \,\Big]
  \;-\;
  \beta \, \KL\!\big( \pol(\cdot \given x) \,\Vert\, \piref(\cdot \given x) \big).
  \label{eq:rlhf}
\end{equation}

The KL term is the regularizer: without it, the policy would collapse onto
whatever response maximizes the (imperfect, learned) reward, exploiting its
errors. $\beta$ sets the trust radius.

% ============================================================
\section{Closed-form optimal policy}
% ============================================================

\begin{proposition}[Optimal policy for the KL-regularized objective]
For a fixed prompt $x$, the policy maximizing \eqref{eq:rlhf} is
\begin{equation}
  \polstar(y \given x)
  = \frac{1}{Z(x)} \, \piref(y \given x) \,
    \exp\!\Big( \tfrac{1}{\beta} \, r(x, y) \Big),
  \qquad
  Z(x) = \sum_{y} \piref(y \given x) \, \exp\!\Big( \tfrac{1}{\beta} r(x, y) \Big).
  \label{eq:optpolicy}
\end{equation}
\end{proposition}

\begin{proof}
Fix $x$ and write the per-prompt objective by expanding the KL:
\[
  J(\pi)
  = \sum_{y} \pi(y \given x)\, r(x,y)
  - \beta \sum_{y} \pi(y \given x) \log \frac{\pi(y \given x)}{\piref(y \given x)}.
\]
Divide by $\beta$ and flip the sign to turn it into a minimization:
\[
  \min_{\pi} \;
  \sum_{y} \pi(y \given x)
  \log \frac{\pi(y \given x)}{\piref(y \given x)}
  - \frac{1}{\beta} \sum_{y} \pi(y \given x)\, r(x,y).
\]
Pull the reward into the log by defining the normalized target
$\tfrac{1}{Z(x)}\piref(y\given x)\exp(r(x,y)/\beta)$ as in \eqref{eq:optpolicy}.
Then the objective equals
\[
  \sum_y \pi(y\given x)
  \log \frac{\pi(y \given x)}
            {\tfrac{1}{Z(x)}\,\piref(y\given x)\exp(r(x,y)/\beta)}
  \;-\; \log Z(x)
  \;=\;
  \KL\!\big(\pi \,\Vert\, \polstar\big) - \log Z(x).
\]
$Z(x)$ does not depend on $\pi$, and KL divergence is minimized (and equals $0$)
exactly when its two arguments match. Hence $\pi = \polstar$, giving
\eqref{eq:optpolicy}.
\end{proof}

\paragraph{Why this matters.} We now know the \emph{shape} of the optimal policy
in closed form. The only obstacle to using it directly is $Z(x)$: summing over
all possible responses $y$ is intractable. DPO's trick is to make $Z(x)$ cancel.

% ============================================================
\section{Reparameterizing the reward}
% ============================================================

Solve \eqref{eq:optpolicy} for the reward by taking logs and rearranging:
\begin{equation}
  r(x, y)
  = \beta \log \frac{\polstar(y \given x)}{\piref(y \given x)}
  + \beta \log Z(x).
  \label{eq:reward-reparam}
\end{equation}

This is the key reparameterization: \emph{any} reward function can be expressed
through its corresponding optimal policy. We do not need a separate reward
network --- the policy \emph{is} an implicit reward model, up to the
prompt-dependent constant $\beta \log Z(x)$.

% ============================================================
\section{Bradley--Terry preferences}
% ============================================================

The Bradley--Terry model says the probability that $\yw$ is preferred over $\yl$
is a sigmoid of their reward difference:
\begin{equation}
  p(\yw \succ \yl \given x)
  = \frac{\exp r(x, \yw)}{\exp r(x, \yw) + \exp r(x, \yl)}
  = \sigmoid\!\big( r(x, \yw) - r(x, \yl) \big).
  \label{eq:bt}
\end{equation}

Crucially this depends only on the reward \emph{difference}.

% ============================================================
\section{The DPO loss}
% ============================================================

Substitute the reparameterized reward \eqref{eq:reward-reparam} into the
Bradley--Terry difference \eqref{eq:bt}. The intractable $\beta \log Z(x)$ term
is identical for $\yw$ and $\yl$ (it depends only on $x$), so it cancels:
\begin{align}
  r(x, \yw) - r(x, \yl)
  &= \beta \log \frac{\polstar(\yw \given x)}{\piref(\yw \given x)}
   - \beta \log \frac{\polstar(\yl \given x)}{\piref(\yl \given x)}
   + \underbrace{\beta \log Z(x) - \beta \log Z(x)}_{=\,0}.
\end{align}

Replacing the unknown optimum $\polstar$ with our trainable $\pol$ and fitting
by maximum likelihood over the preference dataset gives the DPO loss:
\begin{empheq}[box=\fbox]{equation}
  \mathcal{L}_{\mathrm{DPO}}(\theta)
  = - \, \E_{(x, \yw, \yl) \sim \D}
  \left[
    \log \sigmoid\!\left(
      \beta \log \frac{\pol(\yw \given x)}{\piref(\yw \given x)}
      - \beta \log \frac{\pol(\yl \given x)}{\piref(\yl \given x)}
    \right)
  \right].
  \label{eq:dpo}
\end{empheq}

Define the per-response \emph{implicit reward} (the log-ratio, also called the
log-probability margin against the reference):
\[
  \hat r_\theta(x, y) \coloneqq \beta \log \frac{\pol(y \given x)}{\piref(y \given x)}.
\]
Then $\mathcal{L}_{\mathrm{DPO}} = -\,\E\big[\log\sigmoid(\hat r_\theta(x,\yw) - \hat r_\theta(x,\yl))\big]$,
which is just binary logistic regression on the reward margin.

% ============================================================
\section{Gradient and intuition}
% ============================================================

Let $u = \hat r_\theta(x, \yw) - \hat r_\theta(x, \yl)$ be the implicit reward
margin. Using $\frac{d}{du}\log\sigmoid(u) = \sigmoid(-u)$, the gradient is
\begin{equation}
  \nabla_\theta \mathcal{L}_{\mathrm{DPO}}
  = -\,\beta\, \E_{\D}\Big[
      \underbrace{\sigmoid\big(-u\big)}_{\text{weight}}
      \big(
        \nabla_\theta \log \pol(\yw \given x)
        - \nabla_\theta \log \pol(\yl \given x)
      \big)
    \Big].
  \label{eq:grad}
\end{equation}

Reading \eqref{eq:grad}:
\begin{itemize}
  \item The update \textbf{increases} $\log \pol(\yw \given x)$ and
        \textbf{decreases} $\log \pol(\yl \given x)$ --- push up the chosen,
        push down the rejected.
  \item The weight $\sigmoid(-u) = \sigmoid\big(\hat r_\theta(x,\yl) - \hat
        r_\theta(x,\yw)\big)$ is large exactly when the model is \emph{wrong}
        (ranks the rejected response higher) and small once it already ranks
        the pair correctly. The model spends gradient where it is mistaken.
  \item Everything is measured \emph{relative to $\piref$}: the log-ratios prevent
        the policy from trivially inflating both probabilities, and $\beta$
        scales how aggressively margins translate into updates.
\end{itemize}

% ============================================================
\section{What this means for the implementation}
% ============================================================

To compute \eqref{eq:dpo} in code, per batch we need, for each of
$\yw$ and $\yl$:
\begin{enumerate}
  \item $\log \pol(y \given x)$ --- sum of token log-probs under the trainable
        policy (masking out the prompt tokens).
  \item $\log \piref(y \given x)$ --- the same sum under the frozen reference (no
        gradient).
\end{enumerate}
Form the two log-ratios, scale by $\beta$, take the difference, push through
$-\log\sigmoid$, and average. That single scalar is the loss --- no reward
model, no sampling, no RL loop.

\end{document}
